\CatchFileBetweenTags{\AlphaRunningVal}{constants.tex}{AlphaRunningVal}
\CatchFileBetweenTags{\AlphaRunningExperimentalVal}{constants.tex}{AlphaRunningExperimentalVal}
\CatchFileBetweenTags{\AlphaRunningSigma}{constants.tex}{AlphaRunningSigma}

\subsection{Rigidity: The Beta Function (\texorpdfstring{$\beta_0$}{beta\_0})}

The \textbf{Governor} pillar of the Capacity Partition requires structural rigidity: the partition ratios must resist deformation under localized perturbations. The QCD beta function coefficients are the exact geometric encoding of this stiffness.

\paragraph{Geometric Stiffness: The Anti-Screening Term}
The proposed coefficient ``11'' counts the ranks of two geometrically independent contributions:

The Spacetime Anchor ($D \times \chi = 8$) and Symmetry Pressure ($\sigma - \chi = 3$) act on spacetime and internal subspaces respectively. The proposed count adds their ranks. A tensor product of the two identity operators would instead have trace $8\cdot3=24$ (Lean: \texttt{ReviewChecks.tensor\_trace\_check}), so the additive combination is a hypothesis about how the two stiffness contributions compose, not a consequence of their orthogonality. Other integer combinations ($D\times\sigma=20$, $\chi\times\sigma=10$, $D+\chi+\sigma=11$) show that the rule must be fixed independently of the target value.

The proposed vacuum rigidity is then:
\begin{equation}
    \beta_0^{\text{stiff}} = \operatorname{rank}(T_{anchor}) + \operatorname{rank}(T_{pressure}) = 8 + 3 = \mathbf{11}
\end{equation}

This derivation is a \textbf{universality test}: if the geometric construction rules are truly substrate properties rather than group-theoretic accidents, the same rules must apply to all gauge theories. The QED coefficient $4/3$ is not fitted; it is the unique ratio of $2\chi = 4$ (topological load of a virtual pair) to $D-1 = 3$ (spatial dilution). Any other value would require either a different topology (violating $\chi=2$) or different dimensionality (violating $D=4$).

\paragraph{Topological Distribution: The Screening Term}
Fermions screen the charge by distributing the boundary topology ($\chi=2$) across the generation manifold ($n_{gen}=3$, derived strictly from the Triality of the $D_4$ internal sector) which is geometrically isomorphic to the available degrees of freedom in the internal color bulk ($\sigma - \chi = 5 - 2 = 3$). Distributing the boundary evenly across these internal dimensions yields the exact fractional screening per flavor:
\begin{equation}
    \beta_0^{\text{screen}} = \left( \frac{\chi}{n_{gen}} \right) \cdot n_f = \frac{2}{3} n_f
\end{equation}

Combining Rigidity and Screening, we recover the exact QCD Beta Function:
\begin{equation} \label{eq:beta_qcd}
    \beta_0 = (D\chi) + (\sigma-\chi) - \frac{\chi}{n_{gen}} n_f = 11 - \frac{2}{3}n_f
\end{equation}

\textit{Universality check.} The rank composition cannot be the universal rule this section asks for: the weak $SU(2)$ gauge coefficient is $22/3$, and no sum of integer ranks equals it (\texttt{BetaUniversality.rank\_rule\_not\_universal}). The single rule that reproduces both is $(11/3)\,C_A$ with $C_A=\sigma-\chi=3$ for colour and $C_A=\chi=2$ for the weak sector (\texttt{casimir\_form}); the factor $11/3$ is the standard one-loop result (paramagnetic $4$ minus diamagnetic $1/3$ per gauge mode). The coefficients themselves belong with the gauge dynamics of Paper~II, whose Lean states the spin decomposition $(-1)^{2s}((2s)^2-1/3)$ per unit index and derives $b_0=(11/3)N-(2/3)n_f$, hence $9,25/3,23/3,7$ for three to six flavours (\texttt{PMAction.RunningCoefficients.b0\_formula}, \texttt{qcd\_values}); every coefficient used in this paper's running tests is one of these values (\texttt{RunningLinks}). Deriving the two inputs, the spin response $(2s)^2$ and the orbital $1/3$, from the substrate response remains open.

\textbf{Physical Consequence (Asymptotic Freedom):} For $n_f \leq 16$, we have $\beta_0 > 0$, meaning $\alpha_s$ decreases logarithmically at high energy, a phenomenon called \textit{asymptotic freedom}, analogous to \textbf{negative feedback} in control systems. This emerges because Geometric Stiffness ($11$) exceeds Topological Screening ($2n_f/3$) for all Standard Model quark flavors. The strong force weakens at short distances because the vacuum's geometric rigidity dominates the fermion screening contribution.

\CatchFileBetweenTags{\AlphaSTauLinearVal}{constants.tex}{AlphaSTauLinearVal}
\CatchFileBetweenTags{\AlphaSTauThresholdVal}{constants.tex}{AlphaSTauThresholdVal}
\CatchFileBetweenTags{\AlphaSTauThresholdExperimentalVal}{constants.tex}{AlphaSTauThresholdExperimentalVal}
\CatchFileBetweenTags{\AlphaSTauThresholdSigma}{constants.tex}{AlphaSTauThresholdSigma}

\subsubsection{Validation 2: Geometric Scaling (Running to \texorpdfstring{$M_\tau$}{Mtau})}

To test the stiffness coefficient we trace the coupling across energy scales. If each species' response is additive over scale ratios and continuous, it is a coefficient times the logarithm of the scale ratio (Lean: \texttt{ScaleResponse.scale\_homogeneous\_log}); applied to $\alpha_s^{-1}$ this gives the one-loop form
\begin{equation}
\alpha_s^{-1}(\mu) = \alpha_s^{-1}(M) + \frac{\beta_0}{2\pi}\ln\frac{\mu}{M}.
\end{equation}

\textit{Active flavours.} Under the same species-additive response, a flavour screens only above its own mass. Paper~V's quark masses order as $m_c<m_\tau<m_b<M_Z$ (Lean: \texttt{StrongRunning.flavour\_order}), so five flavours are active on $[m_b,M_Z]$ and four on $[m_\tau,m_b]$. Holding $n_f=3$ across the whole interval gives a strictly smaller $\alpha_s^{-1}(m_\tau)$, by the exact excess $\frac{1}{2\pi}\left(\frac43\ln\frac{M_Z}{m_b}+\frac23\ln\frac{m_b}{m_\tau}\right)$ (\texttt{three\_flavours\_inflate}). With the thresholds,
\begin{equation}
\alpha_s^{\text{(1-loop)}}(m_\tau) \approx \AlphaSTauThresholdVal,
\end{equation}
compared with $\AlphaSTauLinearVal$ for three flavours throughout.

\vspace{0.5em}
\begin{tabular}{ll}
\textbf{One-loop, thresholds:} & \AlphaSTauThresholdVal \\
\textbf{Experiment ($\tau$ decays):} & \AlphaSTauThresholdExperimentalVal \\
\textbf{Deviation:}   & $\AlphaSTauThresholdSigma \sigma$ \\
\end{tabular}
\vspace{0.5em}

The one-loop deficit is the expected size of the next order in the coupling: as an external comparison, adding the standard two-loop coefficient with the same thresholds gives $\alpha_s(m_\tau)\approx0.313$. A substrate derivation of the second response coefficient is therefore the next requirement for this test.

\textit{A capacity factor is not supported.} Multiplying the three-flavour value by $(\nu-1)/\nu=15/16$ does not follow from the substrate. A multiplicative factor changes the coupling by the same $1/16$ at every scale, so it cannot be negligible at weak coupling and significant only at strong coupling (\texttt{uniform\_factor\_scale\_free}); applied at the Z mass it would give $\alpha_s<0.1106$ (\texttt{factor\_at\_z}). Normalizing sixteen channel probabilities to one also does not assign the reference channel the share $1/16$ (\texttt{locked\_channel\_share}, \texttt{locked\_share\_not\_forced}). With consistent thresholds the factor would move the estimate further from the data.

\subsubsection{Derivation 3: Geometric Universality --- The QED Beta Coefficient}

Having validated the QCD stiffness coefficients through running coupling predictions, we now demonstrate that these geometric construction rules extend beyond the Strong Force. If the beta function coefficients are truly geometric properties rather than group-theoretic accidents, the same rules must apply universally to all gauge theories. We verify this by deriving the QED coefficient from identical geometric principles.

The QED Beta function coefficient arises from the polarization of the vacuum by virtual pairs:
\begin{enumerate}
    \item \textbf{Topological Load ($2\chi$):} A virtual particle-antiparticle pair creates two topological boundaries ($\chi=2 \times 2 = 4$).
    \item \textbf{Spatial Dilution ($D-1$):} The electric flux distributes over the 3-dimensional spatial volume.
\end{enumerate}

The QED Beta coefficient is the exact ratio of topology to spatial dilution, perfectly matching the standard one-loop result:
\begin{equation}
    \beta_0^{\text{QED}} = \frac{2\chi}{D-1} = \frac{4}{3}
\end{equation}

\subsubsection{Consistency Check: The Z-Pole Running Coupling}

The calculator evaluates
\begin{equation}
\alpha^{-1}(M_Z) = \left[ \alpha^{-1}_{geo}(1-\Delta\alpha) \right] - \Delta^0 \cdot \eta \approx \AlphaRunningVal,
\end{equation}
using the measured on-shell vacuum polarization $\Delta\alpha\approx0.0590$ as an external input.

\vspace{0.5em}
\begin{tabular}{ll}
\textbf{Calculated:} & \AlphaRunningVal \\
\textbf{Experiment ($\overline{\rm MS}$):} & \AlphaRunningExperimentalVal \\
\textbf{Deviation:}   & $\AlphaRunningSigma \sigma$ \\
\end{tabular}
\vspace{0.5em}

\textbf{Scheme caution.} The input $\Delta\alpha$ is the on-shell quantity, whereas the target $127.955$ is the $\overline{\rm MS}$ coupling. The two schemes differ by $\Delta\hat\alpha-\Delta\alpha\approx0.0072$ at leading order, about $1.0$ in $\alpha^{-1}$ (Particle Data Group electroweak review). The approximately unit gap attributed to $\Delta^0\eta$ is therefore of the size of this scheme conversion, and the agreement does not test $\eta$. A test of the overhead needs both sides in one scheme and a substrate-derived $\Delta\alpha$.

\textbf{Not the electroweak interface coupling.} Using this Z-pole value as the coupling of the W and Z readouts would require an interface factor $k=\alpha^{-1}/\alpha^{-1}(M_Z)>1.07$, whereas the measured Z mass requires $1.0362<k<1.0365$ (\texttt{ScaleResponse.zpole\_not\_interface}). Writing $1/k=1-\Delta\alpha+X$ forces $0.0236<X<0.0242$: an anti-screening term opposite in sign to vacuum screening is required (\texttt{antiscreening\_required}).
